% !TeX root = ../教材题库_全部条目_转写工作稿.tex
% Source: XJ-B1-C02-EX-P053-Y0304
% 第2章 / 2.3 一元二次不等式 / 教材 p.53 / PDF 第60页
% OCR未校对：公式与文字均可能有误；未完成数学TeX转写；禁止直接发给学生
% 原题及相邻原解截图（核对用）：../题目原图/XJ-B1-C02-EX-P053-Y0304-01.webp
% 本文件只在转写完成后作为题干使用；OCR 条目尚不是可靠题干。
\begin{bankitem}{2.3 例题 例题3}{第2章 / 2.3 一元二次不等式 / 教材 p.53 / PDF 第60页}{XJ-B1-C02-EX-P053-Y0304}{OCR未校对：公式与文字均可能有误；未完成数学TeX转写；禁止直接发给学生}
\begin{ocrdraft}
[教材p.53 / PDF第60页]
\par
0二0]解不等式一十4x一5>0．
\par
解
\par
（方法一）在不等式两端同时乘一1，可得
\par
一4十5<0．
\par
对于方程一4x十5一0，因为A=（一4）2一4\ensuremath{\times}1\ensuremath{\times}5<0，所以方程没有实数根，
\par
于是函数3，一一4工十5的图象与x轴没有交点，如图2，3一7所示．
\par
由图象得不等式2一行十5<0的解集为，
\par
即不等式亠r2十行一5>0的解集为．
\par
=一4无+5
\par
0
\par
2
\par
3
\par
2
\par
2
\par
y=—x+4．一5
\par
图2．3一7
\par
图2．3一8
\par
（方法二）对于方程亠r2十行一5一0，因为A=42一4\ensuremath{\times}（一1）\ensuremath{\times}（一5）<0，所以
\par
方程没有实数根，
\par
于是函数3，一亠十一5的图象与轴没有交点，如图2．3一8所示．
\par
由图象得不等式一十4一5>0的解集为@．
\end{ocrdraft}
\end{bankitem}
